\chapter{Before submitting} \label{ch:conclusions}

A few checks are worth doing before the final submission.

\begin{itemize}
  \item The cover shows the right programme, the full names of supervisor
    and co-supervisors, the academic year and, if required, the student ID.
  \item The compilation completes without errors, and the log contains no
    \emph{undefined references} or \emph{undefined citations}: run the full
    sequence of Chapter~\ref{ch1} (or \texttt{latexmk}) once more after the
    last change.
  \item The log reports no \emph{overfull boxes}, i.e., no lines exceeding
    the margin; long words or formulas causing them can be fixed by
    rephrasing, or by suggesting hyphenation points with \verb|\-|.
  \item Every figure and table is referred to in the text, and its caption
    can be understood without reading the surrounding text.
  \item Bibliographic entries are complete (authors, title, venue, year) and
    have a DOI whenever one exists.
  \item Placeholder material, such as the examples of this document, the
    dedication or the epigraph if not wanted, has been removed.
\end{itemize}

\section*{Customising the style}

The style is meant to be used as it is, but a few aspects are easy to
change in \texttt{slimthesis.sty}, where they are grouped and commented:
the spacing of the bibliography (\cs{bibitemsep} and the line spacing set by
\cs{AtBeginBibliography}), the colour of links (\texttt{slimblue}), and the
fonts. Keep such changes to a minimum, and note them down: when a new
version of the style is released, they will have to be applied again.

The last page of the thesis acknowledges the SLIM laboratory; it can be
omitted with the \texttt{madeinslim=false} option described in
Section~\ref{sec:programme}.
